%% kapa: exam=1 round=30 year=2019 period=2 form=A subject=accounting subject_ko=회계학 seq=5 number=45
%% source: Q-Net 감정평가사 기출문제 (cst00309) artlSeq=5205205 file=extracted/1008/5205205/2019년 제30회 감정평가사 1차시험 문제지/2019년 제30회 감정평가사 1차시험 문제지/2019년도 제30회 감정평가사 1차 2교시 A형.hwp
%% conversion: kordoc 4.24.51 -> tools/gen1.py (auto); review=proofread-sonnet; flags=none
\begin{question}{45}
무형자산의 회계처리에 관한 설명으로 옳은 것을 모두 고른 것은?
\begin{framedtext}
ㄱ. 내용연수가 비한정적인 무형자산은 상각하지 않고, 무형자산의 손상을 시사하는 징후가 있을 경우에 한하여 손상검사를 수행해야 한다. \par
ㄴ. 무형자산을 창출하기 위한 내부 프로젝트를 연구단계와 개발단계로 구분할 수 없는 경우에는 그 프로젝트에서 발생한 지출은 모두 연구단계에서 발생한 것으로 본다. \par
ㄷ. 브랜드, 제호, 출판표제, 고객목록 및 이와 실질이 유사한 항목은 그것을 외부에서 창출하였는지 또는 내부적으로 창출하였는지에 관계없이 취득이나 완성 후의 지출은 발생시점에 무형자산의 원가로 인식한다. \par
ㄹ. 내용연수가 유한한 무형자산의 잔존가치는 적어도 매 회계연도 말에는 검토하고, 잔존가치의 변동은 회계추정의 변경으로 처리한다. \par
ㅁ. 무형자산은 처분하는 때 또는 사용이나 처분으로부터 미래경제적효익이 기대되지 않을 때 재무상태표에서 제거한다. \par
\end{framedtext}
\begin{choices}
  \item ㄱ, ㄴ, ㄷ
  \item ㄱ, ㄷ, ㄹ
  \item ㄱ, ㄹ, ㅁ
  \item ㄴ, ㄷ, ㅁ
  \item ㄴ, ㄹ, ㅁ
\end{choices}
\end{question}
\answer{5}
